\documentclass[11pt,a4paper]{article}
\usepackage[T1]{fontenc}
\usepackage[utf8]{inputenc}
\usepackage{amsmath,amssymb,amsthm}
\usepackage[margin=1in]{geometry}
\usepackage[colorlinks=true,linkcolor=blue,urlcolor=blue]{hyperref}
\theoremstyle{plain}
\newtheorem{theorem}{Theorem}
\newtheorem{lemma}[theorem]{Lemma}
\newtheorem{proposition}[theorem]{Proposition}
\newtheorem{corollary}[theorem]{Corollary}
\theoremstyle{definition}
\newtheorem{remark}[theorem]{Remark}

\title{The empirical closed form for OEIS A184568, proved}
\author{Adrian Perez Fontelles\\ \small Independent researcher}
\date{1 September 2026}
\begin{document}
\maketitle

\begin{abstract}
OEIS A184568 carries an empirical closed form for $a(n)$ --- not a recurrence relating
consecutive terms but an explicit formula. It is true, and it is decidable rather than
empirical. The entry counts arrays subject to a condition on a bounded window of consecutive lines, so the lines of an array are the
vertices of a finite digraph, an array is a walk in it, and $a(n)$ is a walk count on
$S=9141$ vertices. Any function of the form $\sum_j p_j(n)\lambda_j^{\,n}$ satisfies the
monic linear recurrence whose characteristic polynomial is
$\prod_j(x-\lambda_j)^{\deg p_j+1}$; here that polynomial is $q(x)=(x-1)^{26}$, of degree
$26$. Two things are then checked exactly: that the walk count satisfies the same
recurrence from some index on, which the Cayley--Hamilton theorem reduces to a finite
computation, and that the two sequences agree at $26$ consecutive indices past that
point. A monic recurrence determines every later term from its predecessors, so agreement
there is agreement for good.
\end{abstract}

\noindent\small 2020 Mathematics Subject Classification. 05A15, 11B37, 15A18.\normalsize

\section{The conjecture}

OEIS A184568 is ``Number of (n+2)X5 0..3 arrays with each 3X3 subblock having rows and columns in lexicographically nondecreasing order.''. It has offset $1$ and begins
\[
545662, 6803631, 57374460, 380532059, 2113138210, 10202200416, 43935544294, 171891306894,\ \dots
\]
The entry states, as a conjecture contributed by its author and never marked settled:
\begin{quote}
Empirical: a(n) = (1/19437606570590208000000)*n\^{}25 + (9661/620448401733239439360000)*n\^{}24 + (19289/8617338912961658880000)*n\^{}23 + (647623/1926858390475898880000)*n\^{}22 + (7818869/185785244260761600000)*n\^{}21 + (179204479/50048269882490880000)*n\^{}20 + (6120818579/29194824098119680000)*n\^{}19 + (143209419911/16133981738434560000)*n\^{}18 + (923628988921/3259390250188800000)*n\^{}17 + (7642143794531/1084637427793920000)*n\^{}16 + (243608939835379/1739939207086080000)*n\^{}15 + (1936163562822023/852215121838080000)*n\^{}14 + (304569580619323823/9942509754777600000)*n\^{}13 + (29005828878292128217/83517081940131840000)*n\^{}12 + (233821500453548659/70300574023680000)*n\^{}11 + (7270110226386747353/271159356948480000)*n\^{}10 + (2420495398281155097653/13444984782028800000)*n\^{}9 + (23339250455280961243/23518923816960000)*n\^{}8 + (12862196280962074967/2910169866240000)*n\^{}7 + (66716838431089951113047/4257578514309120000)*n\^{}6 + (473698445981492045742353/10841056402176000000)*n\^{}5 + (4365469476964635381733/46461670295040000)*n\^{}4 + (6119481830014897449583/41557382875008000)*n\^{}3 + (1591545230217707383/10601373182400)*n\^{}2 + (106282373513353/1274816400)*n + 5965
\end{quote}
As of the ``Last modified'' line on the live entry (Oct 03 2025 02:45:47, revision 9) this is still
recorded as empirical, and nothing on the entry records it as proved. Write
\[
f(n)\;=\;\frac{n^{25}}{19437606570590208000000} + \frac{9661 n^{24}}{620448401733239439360000} + \frac{19289 n^{23}}{8617338912961658880000} + \frac{647623 n^{22}}{1926858390475898880000} + \frac{7818869 n^{21}}{185785244260761600000} + \frac{179204479 n^{20}}{50048269882490880000} + \frac{6120818579 n^{19}}{29194824098119680000} + \frac{143209419911 n^{18}}{16133981738434560000} + \frac{923628988921 n^{17}}{3259390250188800000} + \frac{7642143794531 n^{16}}{1084637427793920000} + \frac{243608939835379 n^{15}}{1739939207086080000} + \frac{1936163562822023 n^{14}}{852215121838080000} + \frac{304569580619323823 n^{13}}{9942509754777600000} + \frac{29005828878292128217 n^{12}}{83517081940131840000} + \frac{233821500453548659 n^{11}}{70300574023680000} + \frac{7270110226386747353 n^{10}}{271159356948480000} + \frac{2420495398281155097653 n^{9}}{13444984782028800000} + \frac{23339250455280961243 n^{8}}{23518923816960000} + \frac{12862196280962074967 n^{7}}{2910169866240000} + \frac{66716838431089951113047 n^{6}}{4257578514309120000} + \frac{473698445981492045742353 n^{5}}{10841056402176000000} + \frac{4365469476964635381733 n^{4}}{46461670295040000} + \frac{6119481830014897449583 n^{3}}{41557382875008000} + \frac{1591545230217707383 n^{2}}{10601373182400} + \frac{106282373513353 n}{1274816400} + 5965.
\]

\section{The count is a walk count}

The entry's defining condition is local: it is a condition on a bounded window of consecutive lines. Take as vertices the
admissible configurations of such a window and put an edge from $u$ to $v$ when $v$ may
follow $u$, which the condition decides by inspection. An admissible array is then exactly a
walk, so with $M$ the adjacency matrix of that digraph on $S=9141$ vertices, and $u,w$ the
vectors recording which windows may start and end an array,
\[
a(n)\;=\;u^{\!\top}M^{\,g(n)}w
\]
for the linear function $g$ that the entry's shape supplies. In particular $a$ is
$C$-finite: it satisfies some monic integer linear recurrence, though not necessarily a short
one.

\section{A closed form is a recurrence}

\begin{lemma}\label{lem:ann}
Let $f(m)=\sum_{j}p_j(m)\lambda_j^{\,m}$ with the $\lambda_j$ distinct and the $p_j$
polynomials, and put $q(x)=\prod_j(x-\lambda_j)^{\deg p_j+1}$. Then $q$ applied to the shift
operator annihilates $f$: writing $q(x)=x^{r}-\sum_{i=1}^{r}c_ix^{r-i}$,
\[
f(m)\;=\;\sum_{i=1}^{r}c_i\,f(m-i)\qquad\text{for every }m.
\]
\end{lemma}

\begin{proof}
The shift operator $E$ acts on $m\mapsto\lambda^{m}p(m)$ as
$(E-\lambda)\bigl(\lambda^{m}p(m)\bigr)=\lambda^{m+1}\bigl(p(m+1)-p(m)\bigr)$, and
$p(m+1)-p(m)$ has degree one less than $p$. So $(E-\lambda)^{\deg p+1}$ kills
$\lambda^{m}p(m)$, and the factors of $q$ commute.
\end{proof}

Here $q(x)=(x-1)^{26}$, of degree $26$; the closed form is a polynomial in $n$, so this is $(x-1)^{26}$, and the recurrence it encodes is
\begin{equation}\label{eq:rec}
a(n)\;=\;26\,a(n-1) - 325\,a(n-2) + 2600\,a(n-3) - 14950\,a(n-4) + 65780\,a(n-5) - 230230\,a(n-6) + 657800\,a(n-7) - 1562275\,a(n-8) + 3124550\,a(n-9) - 5311735\,a(n-10) + 7726160\,a(n-11) - 9657700\,a(n-12) + 10400600\,a(n-13) - 9657700\,a(n-14) + 7726160\,a(n-15) - 5311735\,a(n-16) + 3124550\,a(n-17) - 1562275\,a(n-18) + 657800\,a(n-19) - 230230\,a(n-20) + 65780\,a(n-21) - 14950\,a(n-22) + 2600\,a(n-23) - 325\,a(n-24) + 26\,a(n-25) - a(n-26).
\end{equation}

\section{The proof}

\begin{lemma}\label{lem:crit}
With $q$ as above, $a$ satisfies \eqref{eq:rec} for all large $n$ if and only if
$u^{\!\top}M^{\,j}q(M)w=0$ for every $j\ge0$, and it is enough to check $j<S$.
\end{lemma}

\begin{proof}
Substituting the walk expression, the residual $a(n)-\sum_ic_ia(n-i)$ equals
$u^{\!\top}M^{\,g(n)-r}q(M)w$ up to the fixed linear reparametrisation $g$. For the bound,
$M^{S}$ is an integer combination of $I,M,\dots,M^{S-1}$ by Cayley--Hamilton, so the values
$u^{\!\top}M^{j}q(M)w$ satisfy the monic recurrence given by the characteristic polynomial of
$M$; $S$ consecutive zeros therefore force all later ones, and the last nonzero value pins the
threshold exactly.
\end{proof}

\begin{theorem}
$a(n)=f(n)$ for every $n\ge1$.
\end{theorem}

\begin{proof}
The vector $w'=q(M)w$ was formed by $26$ matrix--vector products in exact integer
arithmetic, and $u^{\!\top}M^{j}w'$ evaluated for $j=0,1,\dots$, again exactly. Those integers
vanish from the point corresponding to $n=27$ onwards, and the run was continued until the
working vector was identically zero, which settles every larger $j$ at once. So $a$ satisfies
\eqref{eq:rec} for every $n>26$. By Lemma~\ref{lem:ann} $f$ satisfies \eqref{eq:rec}
identically. Both were then evaluated in exact arithmetic at the $26$ consecutive indices
$n=27,\dots,52$ and agree at each. A monic recurrence of order $26$
determines $a(m)$ from $a(m-1),\dots,a(m-26)$, and for every $m\ge53$ both
$m>26$ and $m-26\ge27$ hold, so induction on $m$ gives $a(m)=f(m)$ for every
$m\ge27$.
Below $n=27$ the recurrence argument gives nothing, since \eqref{eq:rec} is only
established for $a$ from $n=27$ on. The remaining indices $n=1,\dots,26$ were
therefore checked one at a time, directly against the walk count in exact integer arithmetic,
and agree.
\end{proof}

The range is tight in the sense that was checked: the closed form holds from the entry's first term onwards.

\section{Verification}

Four checks, each independent of the algebra above.

First, the walk model reproduces every one of the $16$ terms the entry publishes,
exactly, in integer arithmetic. This is what ties the model to the entry: the model is built
by reading the entry's English, and a misreading gives a different digraph and different
counts.

Second, the closed form was evaluated in exact rational arithmetic --- not floating point ---
at every published term from $n=1$ on, and agrees with the entry's own DATA at each.

Third, the annihilation test of Lemma~\ref{lem:crit} was carried out exactly, and the model
and the closed form were then compared directly at sixty consecutive indices beyond
$n=1$, far past where the proof already settles the matter.

Fourth, the closed form was deliberately perturbed --- by a constant and by a term linear in
$n$ --- and each perturbed form was rejected by the same comparison, so the test is not one
that anything would pass.

\begin{thebibliography}{9}
\bibitem{oeis} The OEIS Foundation, \emph{The On-Line Encyclopedia of Integer
Sequences}, \url{https://oeis.org}, 2026. Sequence A184568.
\bibitem{stanley} R.~P.~Stanley, \emph{Enumerative Combinatorics}, Volume 1, 2nd ed.,
Cambridge University Press, 2011. (Transfer-matrix method, Section 4.7; $C$-finite sequences,
Section 4.1.)
\bibitem{kp} M.~Kauers and P.~Paule, \emph{The Concrete Tetrahedron}, Springer, 2011.
(Closed forms and linear recurrences with constant coefficients.)
\bibitem{horn} R.~A.~Horn and C.~R.~Johnson, \emph{Matrix Analysis}, 2nd ed., Cambridge
University Press, 2013.
\end{thebibliography}

\end{document}
